% macros.tex — the LaTeX Suite chords, as macros.
%
% Each one reproduces the exact output of the Obsidian snippet it is named
% after, so what you draft in the vault and what you type here render the same.
% Capitalised names throughout: the lowercase spellings belong to the physics
% package, which is already loaded and whose \abs, \norm, \eval, \dv, \pdv and
% \qty remain available unchanged.

% --- Function application ----------------------------------------------
% The suite's \!\left(...\right) idiom. These three cover ~60 chords: every
% trig, log, and \operatorname chord differs only in the name it wraps.
\newcommand{\ap}[2]{#1\!\left(#2\right)}            % \ap{\sin}{x}
\newcommand{\apn}[3]{#1^{#2}\!\left(#3\right)}      % Sin2, SinN, LogN, ...
\newcommand{\opn}[2]{\operatorname{#1}\!\left(#2\right)}  % rank, tr, span, ...

% --- Big operators with limits -----------------------------------------
\newcommand{\biglim}[4]{#1\limits_{#2}^{#3}\,#4}
\newcommand{\Sum}[3]{\biglim{\sum}{#1}{#2}{#3}}
\newcommand{\Prd}[3]{\biglim{\prod}{#1}{#2}{#3}}
\newcommand{\BigU}[3]{\biglim{\bigcup}{#1}{#2}{#3}}
\newcommand{\BigI}[3]{\biglim{\bigcap}{#1}{#2}{#3}}
\newcommand{\BigW}[3]{\biglim{\bigwedge}{#1}{#2}{#3}}
\newcommand{\BigV}[3]{\biglim{\bigvee}{#1}{#2}{#3}}

% One limit below, nothing above.
\newcommand{\underlim}[3]{#1\limits_{#2}\,#3}
\newcommand{\Max}[2]{\underlim{\max}{#1}{#2}}
\newcommand{\Min}[2]{\underlim{\min}{#1}{#2}}
\newcommand{\Sup}[2]{\underlim{\sup}{#1}{#2}}
\newcommand{\Inf}[2]{\underlim{\inf}{#1}{#2}}
\newcommand{\Amax}[2]{\underlim{\operatorname*{argmax}}{#1}{#2}}
\newcommand{\Amin}[2]{\underlim{\operatorname*{argmin}}{#1}{#2}}
\newcommand{\Lim}[3]{\lim\limits_{#1 \to #2}\,#3}

% --- Derivatives -------------------------------------------------------
% Upright d, matching the suite; the physics package's \dv uses its own.
\newcommand{\Der}[2]{\frac{\mathrm{d}#1}{\mathrm{d}#2}}
\newcommand{\DerN}[3]{\frac{\mathrm{d}^{#1}#2}{\mathrm{d}#3^{#1}}}
\newcommand{\Dd}[2]{\frac{\mathrm{d}}{\mathrm{d}#1}\,#2}
\newcommand{\DdN}[3]{\frac{\mathrm{d}^{#1}}{\mathrm{d}#2^{#1}}\,#3}
\newcommand{\Pa}[2]{\frac{\partial #1}{\partial #2}}
\newcommand{\PaN}[3]{\frac{\partial^{#1}#2}{\partial #3^{#1}}}
\newcommand{\Pam}[3]{\frac{\partial^2 #1}{\partial #2 \partial #3}}
\newcommand{\Ppa}[2]{\frac{\partial}{\partial #1}\left(#2\right)}
\newcommand{\PpaN}[3]{\frac{\partial^{#1}}{\partial #2^{#1}}\left(#3\right)}

% --- Integrals ---------------------------------------------------------
\newcommand{\Int}[2]{\int #1 \, \mathrm{d}#2}
\newcommand{\DInt}[4]{\int_{#1}^{#2} #3 \, \mathrm{d}#4}
\newcommand{\OInt}[3]{\oint_{#1} #2 \, \mathrm{d}#3}
\newcommand{\IInt}[3]{\iint #1 \, \mathrm{d}#2\,\mathrm{d}#3}
\newcommand{\TInt}[4]{\iiint #1 \, \mathrm{d}#2\,\mathrm{d}#3\,\mathrm{d}#4}
\newcommand{\OInf}[2]{\DInt{0}{\infty}{#1}{#2}}
\newcommand{\Infi}[2]{\DInt{-\infty}{\infty}{#1}{#2}}
\newcommand{\Evb}[3]{\left[ #1 \right]_{#2}^{#3}}

% --- Delimiters --------------------------------------------------------
\newcommand{\Lr}[1]{\left( #1 \right)}
\newcommand{\Avg}[1]{\left\langle #1 \right\rangle}
\newcommand{\Abs}[1]{\left\lvert #1 \right\rvert}
\newcommand{\Norm}[1]{\left\lVert #1 \right\rVert}
\newcommand{\Snorm}[1]{\left\lVert #1 \right\rVert_{2}}
\newcommand{\Fnorm}[1]{\left\lVert #1 \right\rVert_{F}}
\newcommand{\Ceil}[1]{\left\lceil #1 \right\rceil}
\newcommand{\Floor}[1]{\left\lfloor #1 \right\rfloor}
\newcommand{\Set}[1]{\left\{ #1 \right\}}
\newcommand{\Sbd}[2]{\left\{ #1 \mid #2 \right\}}
\newcommand{\Dprd}[2]{\left\langle #1, #2 \right\rangle}
\newcommand{\Lb}[2]{\left[ #1, #2 \right]}
\newcommand{\Sk}[1]{\left[ #1 \right]_{\times}}

% --- Sequences ---------------------------------------------------------
\newcommand{\Sset}[2]{\left\{ #1_1, #1_2, \ldots, #1_{#2} \right\}}
\newcommand{\Sarr}[2]{\left( #1_1, #1_2, \ldots, #1_{#2} \right)}
\newcommand{\Csset}[2]{\left\{ #1_1, \ldots, #1_{#2} \right\}}
\newcommand{\Csarr}[2]{\left( #1_1, \ldots, #1_{#2} \right)}
\newcommand{\Iset}[3]{\left\{ #1_{#2} \right\}_{#2=1}^{#3}}

% --- Physics and quantum mechanics -------------------------------------
\newcommand{\kbt}{k_{B}T}
\newcommand{\msun}{M_{\odot}}
\newcommand{\Iso}[3]{{}^{#1}_{#2}\mathrm{#3}}
\newcommand{\Outer}[1]{\ket{#1}\bra{#1}}
\newcommand{\Brk}[2]{\braket{#1 | #2}}
\newcommand{\Cis}[1]{\cos\!\left(#1\right) + \mathrm{i}\sin\!\left(#1\right)}
